\subsection{What the three auxiliary units actually annihilate}
\label{sec:kt-auxiliary-unit-bridge}

\begin{theorem}[KT-AUXILIARY-UNIT-SELMER-BRIDGE]
Let $p>3$ be prime, let $D<0$ be a fundamental discriminant
prime to $p$, and put
\[
 E=\mathbf Q(\sqrt D),\quad L=E(\mu_p),\quad
 \chi=\chi_D,\quad A=\operatorname{Cl}(L)\{p\}.
                                                               \tag{AB.1}
\]
Let $a\ge0$ be even, with either $a=0$ or
$a\not\equiv0\pmod{p-1}$.  If
\[
                         B_{1,\chi\omega^a}\in\mathbf Z_p^\times,
                                                               \tag{AB.2}
\]
then
\[
             A^{\chi\omega^{-a}}=0,
             \qquad A^{\chi\omega^{a+1}}=0.               \tag{AB.3}
\]
Here a superscript denotes the eigenspace for the indicated
Galois character, with the fixed $p$-adic Teichm\"uller embedding.

Write $\varepsilon_p$ for the $p$-adic cyclotomic character and
$T_a=\mathbf Q_p/\mathbf Z_p(\chi\varepsilon_p^{-a})$.
Define the strict group by its actual local conditions:
\[
 \mathcal S_p(T_a)=\ker\left(H^1(\mathbf Q,T_a)\longrightarrow
 H^1(\mathbf Q_p,T_a)\oplus
       \prod_{q\ne p}H^1(I_q,T_a)\right),                 \tag{AB.4}
\]
where $q$ ranges over finite rational primes.  Then
\[
                            \mathcal S_p(T_a)=0.          \tag{AB.5}
\]
The same assertion holds with $T_a$ replaced by
$\mathbf F_p(\chi\omega^{-a})$.

For $a=0,e,2e$ under $e\ge2$ even and
$(p-1)\nmid2e$, the three units therefore give the following
exact character correspondence:
\[
\begin{array}{c|c|c|c}
\text{unit}&\text{odd class component}&\text{even mirror}
  &\text{strict coefficient module}\\ \hline
B_{1,\chi}&\chi&\chi\omega&\mathbf Q_p/\mathbf Z_p(\chi)\\
B_{1,\chi\omega^e}&\chi\omega^{-e}&\chi\omega^{e+1}
 &\mathbf Q_p/\mathbf Z_p(\chi\varepsilon_p^{-e})\\
B_{1,\chi\omega^{2e}}&\chi\omega^{-2e}&\chi\omega^{2e+1}
 &\mathbf Q_p/\mathbf Z_p(\chi\varepsilon_p^{-2e})
\end{array}                                                   \tag{AB.6}
\]
Every component in the second and third columns is negative
under $\operatorname{Gal}(L/\mathbf Q(\mu_p))$.
In contrast, the original rational component $\omega^{-e}$
is positive under that group.  Thus (AB.3)--(AB.6) do not,
by themselves, assert vanishing of the original rational
strict group with coefficient $\mathbf Q_p/\mathbf Z_p(-e)$.
\end{theorem}

\begin{proof}
Set $f=|D|$ and $G=\operatorname{Gal}(L/\mathbf Q)$.
The extensions $E$ and $\mathbf Q(\mu_p)$ are disjoint,
so $|G|=2(p-1)$ is prime to $p$.  All character values
under consideration lie in $\mathbf Z_p$, and integral
character projections are available.

First suppose $a\not\equiv0\pmod{p-1}$.  Then
$\eta=\chi\omega^a$ is primitive of conductor $fp$.
Use the Stickelberger element for the abelian field $L$,
with $\sigma_b$ induced by $\zeta_{fp}\mapsto\zeta_{fp}^b$:
\[
 \Theta_L=\frac1{fp}\sum_{\substack{1\le b\le fp\\(b,fp)=1}}
                     b\sigma_b^{-1}\in\mathbf Q[G].       \tag{AB.7}
\]
For every positive integer $c$ prime to $fp$,
$(c-\sigma_c)\Theta_L$ is integral and annihilates
$\operatorname{Cl}(L)$ by Stickelberger's theorem for an
abelian subfield, not merely for a full cyclotomic field.
Integrality can also be seen before restriction to $L$:
the coefficient at $\sigma_b^{-1}$ is
$(cb-[cb]_{fp})/(fp)\in\mathbf Z$.

On the $\eta^{-1}$ eigenspace, this annihilator acts by
\[
                  (c-\eta^{-1}(c))B_{1,\eta}.             \tag{AB.8}
\]
Indeed evaluation of $\sigma_b^{-1}$ at $\eta^{-1}$
gives $\eta(b)$, so the normalization and inverse in
(AB.8) follow directly from the defining Bernoulli sum.
Choose $c\equiv1\pmod p$ and $\chi(c)=-1$, with
$(c,f)=1$, by the Chinese remainder theorem.  Then
$c-\eta^{-1}(c)\equiv2\pmod p$.  Under (AB.2), (AB.8)
is a unit annihilator; hence $A^{\chi\omega^{-a}}=0$.
There is no assumption $p\nmid\varphi(f)$.

For $a=0$, the imaginary quadratic class-number formula
shows that $p\nmid h(E)$.  With
$H=\operatorname{Gal}(L/E)$, extension after norm is
$\sum_{h\in H}h$.  It acts as the unit $p-1$ on $A^H$,
while norm has target $\operatorname{Cl}(E)\{p\}=0$.
Thus $A^H=0$, proving $A^\chi=0$.  This treatment avoids
inserting the imprimitive Euler factor $1-\chi(p)$, which
would vanish when $p$ splits in $E$.

For the mirror assertion, write
$\alpha=\chi\omega^{a+1}$ and $\xi=\chi\omega^{-a}$.
Then $\alpha$ is even and $\xi=\omega\alpha^{-1}$ is
odd and different from $\omega$.  Apply the full
Kummer-radical argument of the carried K17-REFLECTOR to
$L$.  More explicitly, the quotient
$A^\alpha/pA^\alpha$ gives an everywhere unramified
elementary abelian $p$-extension of $L$, stable under $G$.
Its Kummer radical $V\subset L^\times/L^{\times p}$ has
character $\xi$ by the equivariant pairing into $\mu_p$.
Unramifiedness implies that $(x)=\mathfrak a^p$ for every
representative $x$ of $V$.  Sending $x$ to
$[\mathfrak a]\in A[p]$ is well defined and $G$-equivariant;
its kernel lies in the unit group modulo $p$th powers.
For a CM field, that unit group's odd part consists only
of roots of unity: $u/\overline u$ is a root of unity,
and $u^2=(u/\overline u)(u\overline u)$.
Here the $p$-primary roots of unity are exactly $\mu_p$
because $[L:\mathbf Q]=2(p-1)$ is prime to $p$.
Their character is $\omega$, different from $\xi$.
Consequently $V\hookrightarrow A[p]^\xi$ is injective.
The first part makes its target zero, so
$A^\alpha/pA^\alpha=0$.  Finiteness gives $A^\alpha=0$.

Put $W_a=T_a[p]=\mathbf F_p(\chi\omega^{-a})$.
Since $p\nmid|G|$, restriction identifies global
$H^1(\mathbf Q,W_a)$ with the $G$-invariant part of
$H^1(L,W_a)$.  A class in the residual strict group
restricts to an everywhere unramified homomorphism on
$G_L$: away from $p$ its inertia restriction is zero,
and at $p$ its entire local class is zero.  Global class
field theory therefore embeds that strict group into
\[
                 \operatorname{Hom}_G(A/pA,W_a).          \tag{AB.9}
\]
This is zero because the first component in (AB.3) is zero.
The equivariance convention is explicit: a homomorphism
$f$ in (AB.9) obeys $f(gx)=\chi(g)\omega(g)^{-a}f(x)$;
it sees $A^{\chi\omega^{-a}}$, not its inverse.

It remains to pass from residual to divisible coefficients
without changing the local conditions.  Complex conjugation
acts as $-1$ on $T_a$, so $H^0(\mathbf Q,T_a)=0$ and
\[
              H^1(\mathbf Q,W_a)\simeq
                       H^1(\mathbf Q,T_a)[p].            \tag{AB.10}
\]
For $q\ne p$, the cyclotomic character is unramified, and
inertia acts through the order-at-most-two character $\chi$.
Hence $T_a^{I_q}$ is either $0$ or $T_a$, and is
$p$-divisible.  At $p$, if $a\not\equiv0\pmod{p-1}$,
some inertia element acts nontrivially modulo $p$, so
$H^0(\mathbf Q_p,T_a)=0$.  If $a=0$, this group is
either $0$ or $T_a$, according as $\chi(p)=-1$ or $1$;
it is again $p$-divisible.  The long exact sequences for
$0\to W_a\to T_a\xrightarrow{p}T_a\to0$ consequently
make every residual local $H^1$ map in (AB.4) injective.
Thus (AB.10) identifies the $p$-torsion in
$\mathcal S_p(T_a)$ with the residual strict group,
which is zero by (AB.9).  A nonzero $p$-primary torsion
group has a nonzero element of order $p$, proving (AB.5).

Finally let $\tau$ be the nontrivial automorphism of
$L/\mathbf Q(\mu_p)$.  It has $\chi(\tau)=-1$ and
$\omega(\tau)=1$.  This proves every sign assertion
in (AB.6) and the stated distinction from the rational
character $\omega^{-e}$.
\end{proof}

\paragraph{Imported theorem and exact proof scope.}
The annihilation input is L.~C.~Washington,
\emph{Introduction to Cyclotomic Fields}, first edition,
Springer (1982), \S6.2: definition on p.~93,
Lemma~6.9 and Theorem~6.10 on p.~94, with descent from
the full cyclotomic field to the arbitrary abelian field
on pp.~99--100.  The accessible book OCR was read for
the definition, lemma, theorem, and concluding descent;
the preceding Gauss-sum factorization proof is imported,
not claimed freshly checked.  The hosting page labels
the upload as the 1997 edition, but its actual title
and copyright pages show the 1982 first edition.
The publisher's second-edition chapter is
\url{https://doi.org/10.1007/978-1-4612-1934-7_6};
its abstract independently states the arbitrary-abelian-field
scope, while its paywalled chapter text was not retrieved.
The mirror proof above is a fully spelled-out adaptation
of the already checked K17-REFLECTOR, carried
\texttt{K17.tex}, lines 4533--4612; that argument cites
Sharifi's Theorem~4.4.2 and Corollary~4.4.4, pp.~111--113.
No new claim identifies (AB.4) with a different deformation
problem's local conditions.  Such an identification must
be checked against that problem's actual local predicates.

\begin{lemma}[KT-RATIONAL-POSITIVE-ROOT-UNIT]
Let $p>3$, let $b\ge2$ be even, and suppose
$(p-1)\nmid b$.  If $B_b/b$ is a $p$-adic unit, then
\[
              H^1(G_{\mathbf Q,\{p\}},
                         \mathbf F_p(\omega^b))=0.       \tag{AB.11}
\]
Here $G_{\mathbf Q,\{p\}}$ denotes the Galois group of the
maximal extension unramified at finite primes other than
$p$; no strict condition at $p$ is imposed in (AB.11).
\end{lemma}
\begin{proof}
Let $b_0\in\{2,4,\ldots,p-3\}$ be congruent to $b$
modulo $p-1$.  The ordinary Kummer congruence, equivalently
Sun's Lemma~8.1(a) with trivial character and difference
order one, gives $B_b/b\equiv B_{b_0}/b_0\pmod p$;
all intervening Euler factors are congruent to one.
Herbrand's implication in the carried K17-REFLECTOR proof
therefore gives
\[
 \operatorname{Cl}(\mathbf Q(\mu_p))\{p\}^{\omega^{1-b}}=0.
                                                               \tag{AB.12}
\]
Set $F=\mathbf Q(\mu_p)$, $\Delta=\operatorname{Gal}(F/\mathbf Q)$,
and $S$ equal to its unique prime over $p$.  The $S$-unit
Kummer exact sequence is
\[
 0\longrightarrow\mathcal O_F[1/p]^\times/
                    \mathcal O_F[1/p]^{\times p}
 \longrightarrow H^1(G_{F,S},\mu_p)
 \longrightarrow\operatorname{Cl}(\mathcal O_F[1/p])[p]
 \longrightarrow0.                                      \tag{AB.13}
\]
The prime in $S$ is principal, so the last group equals
$\operatorname{Cl}(F)[p]$.  The only odd character in
ordinary units modulo $p$th powers is $\omega$.
Inverting the unique prime adds only a trivial-character
quotient, since its valuation is $\Delta$-invariant.
Thus the $\omega^{1-b}$ component of the first term is
zero: this character is odd and differs from $\omega$.
The corresponding component of the last term is zero
by (AB.12).  It follows that the middle component is zero.
Finally twist (AB.13) by $\omega^{b-1}$ and take exact
$\Delta$-invariants.  Restriction identifies the resulting
middle term with (AB.11), because $|\Delta|=p-1$ is a unit.
This proves the claim at every positive weight $b$ in the
stated residue class.
\end{proof}

This additional rational statement removes the positive
coefficient $\omega^b$ under its stated Bernoulli-unit
condition.  It does not remove $\omega^{-b}$ without
the corresponding reflected Bernoulli input.  Whether
these groups govern every arrow in a proposed deformation
graph remains a separate question.
